% !TEX root = ../main.tex
% Primary sources: ATLAS of Finite Groups for the M_{24} class table;
% Dijkgraaf, Pasquier and Roche for transgression; Brown for cyclic
% group cohomology; Gaberdiel, Persson and Volpato and
% Cheng, Duncan and Harvey for Mathieu multiplier systems.

\providecommand{\CM}{C_{M_{24}}}
\providecommand{\Z}{\mathbb{Z}}
\providecommand{\U}{\mathrm{U}\!\left(1\right)}
\providecommand{\res}{\operatorname{res}}
\providecommand{\Sym}{\operatorname{Sym}}

\section*{Anomalous Mathieu classes and cyclic transgression}

\paragraph{The three cohomology groups.}
For \(g\in M_{24}\) simultaneous twining keeps three cohomology groups
separate:
\[
H^2(\CM(g),\U),\qquad
H^3(M_{24},\U),\qquad
H^3(\langle g\rangle,\U).
\]
The first records intrinsic projective representations of the
centraliser.  The second is the Dijkgraaf-Witten class whose slant product
gives a commuting-pair \(2\)-cocycle.  The third is the cyclic restriction
detected by the cusp multiplier.  A class in one of these groups supplies
no class in the other two without an explicit restriction or transgression
map.

\paragraph{Centraliser and cyclic data.}
The ATLAS class table gives the centraliser orders
\[
|\CM(7A)|=|\CM(7B)|=42,\qquad
|\CM(11A)|=11,\qquad
|\CM(23A)|=|\CM(23B)|=23.
\]
For the order-seven classes, \(\langle g\rangle\) is central in
\(\CM(g)\) and has coprime index \(6\).  Schur-Zassenhaus splits the
extension.  The ATLAS power maps contain the order \(14\) and order \(21\)
power-up classes and contain no order \(42\) class, so the complement is
\(S_3\), giving \(\CM(g)\cong C_7\times S_3\).  For the order-eleven and
order-twenty-three classes the centraliser is cyclic.  Thus:
\[
\begin{array}{cccc}
\toprule
g & \CM(g) & H^2(\CM(g),\U) & H^3(\langle g\rangle,\U) \\
\midrule
7A,7B & C_7\times S_3 & 0 & \Z/7 \\
11A & C_{11} & 0 & \Z/11 \\
23A,23B & C_{23} & 0 & \Z/23 \\
\bottomrule
\end{array}
\]

\paragraph{Vanishing of the centraliser Schur multiplier.}
For a finite group \(G\), the exponential sequence
\[
0\longrightarrow \Z \longrightarrow \mathbb{R}\longrightarrow \U
\longrightarrow 0
\]
and \(H^{>0}(G,\mathbb{R})=0\) give
\[
H^n(G,\U)\cong H^{n+1}(G,\Z),\qquad n\geq 1.
\]
For \(C_N\), integral cohomology is
\[
H^{2r}(C_N,\Z)\cong \Z/N,\qquad
H^{2r+1}(C_N,\Z)=0\qquad (r\geq 1).
\]
Thus \(H^2(C_N,\U)=0\) and \(H^3(C_N,\U)\cong \Z/N\).
For \(C_7\times S_3\), Schur's product formula gives
\[
M(C_7\times S_3)
 \cong M(C_7)\oplus M(S_3)\oplus
 (C_7^{\mathrm{ab}}\otimes S_3^{\mathrm{ab}}).
\]
Here \(M(C_7)=0\), \(M(S_3)=0\), and
\(\Z/7\otimes S_3^{\mathrm{ab}}=0\) because
\(S_3^{\mathrm{ab}}\cong\Z/2\).  Hence \(H^2(C_7\times S_3,\U)=0\).
The intrinsic Schur multiplier of \(\CM(g)\) therefore vanishes for all
five classes above.

\paragraph{Dijkgraaf, Pasquier and Roche transgression.}
Let \(\alpha\in Z^3(M_{24},\U)\) be a normalized \(3\)-cocycle.  For
\(x,y\in\CM(g)\), its DPR transgression is
\[
\tau_g(\alpha)(x,y)
 =
 \frac{\alpha(g,x,y)\,\alpha(x,y,g)}{\alpha(x,g,y)}.
\]
This is a \(2\)-cocycle on \(\CM(g)\).  Since \(H^2(\CM(g),\U)=0\) for the
five classes, every \(\tau_g(\alpha)\) is cohomologous to zero on the
centraliser:
\[
[\tau_g(\alpha)]=0\in H^2(\CM(g),\U).
\]
The anomalous cusp multiplier therefore has its cohomological receptacle
in the cyclic restriction channel.  The intrinsic centraliser Schur
multiplier is zero.

\paragraph{Cyclic restriction.}
The group-cohomological datum that remains available from a chosen
\(3\)-cocycle is its restriction to the cyclic subgroup:
\[
\res^{M_{24}}_{\langle g\rangle}[\alpha]
 =
a_g\,\epsilon_N
\in H^3(C_N,\U)\cong \Z/N,
\qquad N=\operatorname{ord}(g),
\]
where \(\epsilon_N\) is the cyclic generator and \(a_g\in\Z/N\).  When a
twining multiplier is identified with this cyclic residue, the multiplier
character at the canonical cusp is
\[
\chi_g(T)=\exp(2\pi i\,a_g/N).
\]
A non-zero \(a_g\) is then the cyclic group-cohomology witness of the
anomaly.  The numerical residue \(a_g\) is multiplier-system data from
the twining genus and is independent of \(H^2(\CM(g),\U)\).  A fraction
such as \(3/7\), \(2/11\), or \(1/23\) becomes a cohomology statement
only after the cusp generator, multiplier convention, and Galois
labelling of \(7A/7B\) and \(23A/23B\) have been fixed.  The calculation
above proves the centraliser vanishing and the conditional cyclic
formalism.  A non-zero cyclic residue requires a specified global
\(M_{24}\) class and multiplier convention.

\paragraph{Geometric scope.}
Geometric order and anomaly residue are distinct tests.  Order \(7\)
belongs to the Nikulin-Mukai symplectic list; orders \(11\) and \(23\)
lie outside it.  The cohomological statement proved here is the
centraliser Schur-multiplier vanishing together with the cyclic
\(H^3\)-receptacle for multiplier residues.
A K3 conformal-field-theory or Mukai-lattice realisation claim requires
separate class-level geometry.

\paragraph{Separation from the \(K3\times E\) Borcherds constants.}
The cyclic anomaly is independent of the Borcherds weight normalisation.
On the CHL ladder
\[
\begin{array}{cccccc}
\toprule
N & 1 & 2 & 3 & 4 & 6\\
\midrule
c_N(0) & 10 & 8 & 6 & 4 & 2\\
\kappa_{\mathrm{BKM}}(\Phi_N) & 5 & 4 & 3 & 2 & 1
\bottomrule
\end{array}
\]
and in every column
\[
\kappa_{\mathrm{BKM}}(\Phi_N)=\frac{c_N(0)}{2}.
\]
For \(N=1\),
\[
D_5=64^{-1}\Delta_5,\qquad
\Phi_1=\Delta_5,\qquad
c_1(0)=10,\qquad
\kappa_{\mathrm{BKM}}(\Delta_5)=5.
\]
For \(K3\times E\),
\begin{center}
\begin{tabular}{ccl}
\toprule
invariant & value & source\\
\midrule
\(\kappa_{\mathrm{cat}}(K3\times E)\) & \(0\) &
\(\chi(\mathcal O_{K3})\chi(\mathcal O_E)=2\cdot 0\)\\
\(\kappa_{\mathrm{ch}}(K3\times E)\) & \(0\) &
\(\sum_q{(-1)}^q h^{0,q}(K3\times E)=1-1+1-1\)\\
\(\kappa_{\mathrm{ch}}^{\mathrm{Heis}}(K3\times E)\) & \(3\) &
Heisenberg-Mukai specialisation\\
\(\kappa_{\mathrm{BKM}}(\mathfrak g_{\Delta_5})\) & \(5\) &
\(c_{\phi_{0,1}}(0,0)/2=10/2\)\\
\(\kappa_{\mathrm{fiber}}(K3)\) & \(24\) &
\(\chi_{\mathrm{top}}(K3)\)\\
\bottomrule
\end{tabular}
\end{center}
At \(d=3\), the native \(\Phi_3\) output is \(E_1\)-chiral on the chosen
reference curve.  The available \(E_2\)-level belongs to a constructed
centre, double, module category, or enhancement datum.  The
anomalous-class calculation above uses only group cohomology and
multiplier restrictions.

\paragraph{Compact hCS boundary conditions.}
A compact or global holomorphic Chern-Simons boundary statement contains
the following data:
\[
\begin{array}{l}
\text{quantum master equation},\\
\text{anomaly cancellation},\\
\text{orientation and framing},\\
\text{descent, sewing and TCFT compatibility},\\
\text{local counterterms},\\
\text{analytic or Novikov completion}.
\end{array}
\]
In complex dimension \(3\), the local BV gauge anomaly is governed by
the quartic invariant
\[
I_4^\rho\in {\left(\Sym^4 \mathfrak g^\vee\right)}^{\mathfrak g}.
\]
Thus a compact \(K3\times E\) hCS realisation of the \(\Delta_5\)
boundary algebra is an additional quantisation theorem, separate from
the finite-group cohomology calculation and from the Borcherds product
identity \(D_5=64^{-1}\Delta_5\).
